% !TeX root = ../../Book3.tex
% 纲要标题：### 23.5*　扫掠、宽度与非极小临界点
% 纲要位置：几何学计划纲要.md，第 3107 行。

% 纲要内容开始
% 
% > 写作范围与外部输入：规划任务：在第20—22章的弱极限工具与第23.4节变分流形之后引入几何 min-max。完整证明黎曼二球面上 Birkhoff 非平凡闭测地线的代表模块，交代离散缩短、宽度为正和极限不塌缩为常值；这不是第28章的平均曲率流。再在闭环境维数3—7及指定系数、扫掠空间和收敛拓扑下，准确外引所选 Almgren–Pitts 极小超曲面存在性，分别说明重数、正则性和指标。有限维山路模型只作动机，不替代几何存在证明；第16.4节谱的极小极大原则与本节另行辨析。
% 
% * 扫掠族、同伦约束与宽度
% * 分段测地线的缩短与 Birkhoff 模型
% * 二球面闭测地线存在性的完整证明
% * 超曲面扫掠与驻值变分流形
% * Almgren–Pitts 存在性的准确版本
% * 维数、重数、正则性与指标的区别
% 
% ---
% 
% 纲要内容结束
% 正文状态：未撰写；本文件目前是框架。
% 规划任务：在第20—22章的弱极限工具与第23.4节变分流形之后引入几何 min-max。完整证明黎曼二球面上 Birkhoff 非平凡闭测地线的代表模块，交代离散缩短、宽度为正和极限不塌缩为常值；这不是第28章的平均曲率流。再在闭环境维数3—7及指定系数、扫掠空间和收敛拓扑下，准确外引所选 Almgren–Pitts 极小超曲面存在性，分别说明重数、正则性和指标。有限维山路模型只作动机，不替代几何存在证明；第16.4节谱的极小极大原则与本节另行辨析。

\OptionalSection{扫掠、宽度与非极小临界点}{g3:sec:2305}

待撰写

\subsection{扫掠族、同伦约束与宽度}
\label{g3:subsec:2305:01}
待撰写

\subsection{分段测地线的缩短与 Birkhoff 模型}
\label{g3:subsec:2305:02}
待撰写

\subsection{二球面闭测地线存在性的完整证明}
\label{g3:subsec:2305:03}
待撰写

\subsection{超曲面扫掠与驻值变分流形}
\label{g3:subsec:2305:04}
待撰写

\subsection{Almgren–Pitts 存在性的准确版本}
\label{g3:subsec:2305:05}
待撰写

\subsection{维数、重数、正则性与指标的区别}
\label{g3:subsec:2305:06}
待撰写
